% Source: Weinberg GR, physical PDF pp. 482--487 (printed pp. 459--464).
% Coverage: Section 14.8, The Steady State Cosmology; equations
%           (14.8.1)--(14.8.23), all intervening unnumbered displays, and
%           the concluding paragraph before the chapter bibliography.
% Figures/tables/footnotes: linked reference markers 57, 63, and 65--67;
%                          no figures, tables, or footnotes.
% Status: source-reviewed and compile-clean; visually proofed.
% Uncertainties: the signature, scale-factor, time-index, divergence, unit,
%                spelling, and corrected cumulative-count notation are
%                documented below; no unresolved source uncertainty remains.
%
% MODERNIZATION CHECK: the printed scale-factor symbol is written a(t);
% time is index 0; the metric (14.8.5) is converted to (-+++); covariant
% divergences are explicit; the obsolete spelling of ``deceleration'' is
% regularized; and units are upright with inverse powers.
% VERIFY SOURCE: physical p. 485 says that an observed value of N(>z) is
% smaller than the prediction (14.8.17), but both the definition and the
% logic require N(<z). The evident inequality slip is corrected in prose.
% LITERAL-OPERATOR AUDIT: the minus signs in (14.8.2), (14.8.3),
% (14.8.16), and (14.8.20)--(14.8.23), and the nested roots in
% (14.8.14) and (14.8.18), were checked directly against the source.

\subsection{The Steady State Cosmology}\label{sec:14.8}

Our work so far has been based on the ``Cosmological Principle'' that the
universe is spatially isotropic and homogeneous. Hermann Bondi and Thomas
Gold\hyperref[ch14-ref-65]{\textsuperscript{65}} have gone one step further,
and have suggested that the universe obeys a ``perfect Cosmological
Principle,'' that it looks the same not only at all points and in all
directions, but at all times. This assumption leads to a steady state model
of the universe, which was suggested at about the same time by Fred
Hoyle,\hyperref[ch14-ref-66]{\textsuperscript{66}} on the basis of an
alteration in the structure of the energy--momentum tensor appearing in
Einstein's field equations. We shall follow the Bondi--Gold approach here,
as more suited to the spirit of the present chapter, and will come back to
Hoyle's theory in the last chapter.

The work of Section~\ref{sec:14.6} shows that the Hubble ``constant''
\(\dot a(t_0)/a(t_0)\) is an observable parameter, so that it must be
independent of the present time \(t_0\) in a steady state model. Letting
\(H\) denote the permanent value of the Hubble constant, we have then
\[
  \frac{\dot a(t)}{a(t)}=H
  \qquad\text{for all \(t\),}
\]
and therefore
\begin{equation}
  a(t)=a(t_0)\exp[H(t-t_0)].
  \tag{14.8.1}\label{eq:14.8.1}
\end{equation}
In this model the deceleration parameter takes the permanent value
\begin{equation}
  q=-\frac{a\ddot a}{\dot a^2}=-1.
  \tag{14.8.2}\label{eq:14.8.2}
\end{equation}
To determine \(k\), we return to the general relation
\eqref{eq:14.6.22} between \(a(t)\) and the luminosity-distance versus
red-shift function \(d_L(z)\), which now reads
\begin{align}
  t_0-t
  =
  \int_0^{\exp[H(t_0-t)]-1}
  &\frac{1}{1+z}
  \left[
    1-k\,a^{-2}(t_0)(1+z)^{-2}d_L^2(z)
  \right]^{-1/2}
  \notag\\[-2pt]
  &{}\times
  \frac{\dd}{\dd z}
  \left[
    (1+z)^{-1}d_L(z)
  \right]\dd z.
  \tag{14.8.3}\label{eq:14.8.3}
\end{align}
Since \(d_L(z)\) is observable, it must now be independent of \(t_0\).
Hence, in order that the integral depend only on \(t-t_0\), not \(t\) or
\(t_0\) separately, it is necessary that
\begin{equation}
  k=0.
  \tag{14.8.4}\label{eq:14.8.4}
\end{equation}
The metric is then
\begin{equation}
  \dd s^2
  =
  -\dd t^2
  +
  a^2(t_0)e^{2H(t-t_0)}
  \left(
    \dd r^2+r^2\dd\theta^2+r^2\sin^2\theta\,\dd\phi^2
  \right).
  \tag{14.8.5}\label{eq:14.8.5}
\end{equation}

This derivation might be challenged on the grounds that the metric
\eqref{eq:14.8.5} was obtained as a special case of the
Robertson--Walker metric, which was derived in Sections~\ref{sec:14.1} and
\ref{sec:14.2} on the basis of a definition of cosmic time that makes no
sense in an unevolving universe. We could avoid this difficulty by viewing
Eq.~\eqref{eq:14.8.5} as the limiting metric for a universe with an
extremely slow rate of evolution. A more satisfactory approach is to derive
Eq.~\eqref{eq:14.8.5} directly from the assumption that the whole
four-dimensional spacetime is maximally symmetric. This assumption is shown
in Section~13.3 to lead to the metric \eqref{eq:13.3.41}, which is identical
to Eq.~\eqref{eq:14.8.5}, except that a factor
\(a(t_0)\exp(-Ht_0)\) must be absorbed into the radial coordinate \(r\).
Comparing Eqs.~\eqref{eq:13.3.41} and \eqref{eq:14.8.5}, we note that the
curvature constant of the four-dimensional spacetime of the steady state
model is
\begin{equation}
  K=H^2.
  \tag{14.8.6}\label{eq:14.8.6}
\end{equation}
The \emph{spacetime} is curved, even though the \emph{space} is flat.

The most remarkable feature of the steady state cosmology is not its
spacetime metric, but rather the necessity of continuous creation of matter.
According to Eq.~\eqref{eq:14.2.21}, the proper distance between any two
comoving galaxies increases as \(a(t)\), so if the average number of
galaxies per unit proper volume is to remain constant, new galaxies must
appear to fill up the holes in the widening comoving coordinate mesh. To put
this formally, we recall that in the comoving coordinate system
\((t,r,\theta,\phi)\), the current vector of the galaxies and the total
energy--momentum tensor are given by Eqs.~\eqref{eq:14.2.11}--%
\eqref{eq:14.2.14} as
\[
  J_G^\mu=n_GU^\mu,
\]
\[
  T^{\mu\nu}
  =
  (\rho+p)U^\mu U^\nu+pg^{\mu\nu},
\]
with
\[
  U^0=1,
  \qquad
  U^r=U^\theta=U^\phi=0.
\]
In accordance with the spirit of the steady state model, we now take
\(n_G,\rho,\) and \(p\) to be constant in time as well as space. Then
\(J_G^\mu\) and \(T^{\mu\nu}\) are not conserved, but rather
\begin{equation}
  \nabla_\mu J_G^\mu
  =
  a^{-3}(t)\frac{\dd}{\dd t}
  \left[
    a^3(t)J_G^0
  \right]
  =
  3n_GH,
  \tag{14.8.7}\label{eq:14.8.7}
\end{equation}
\begin{equation}
  \nabla_\nu T^{0\nu}
  =
  a^{-3}(t)\frac{\dd}{\dd t}
  \left[
    a^3(t)(\rho+p)
  \right]
  =
  3(\rho+p)H.
  \tag{14.8.8}\label{eq:14.8.8}
\end{equation}
That is, a comoving observer using a locally inertial coordinate system will
see galaxies created at a rate \(3H\) per existing galaxy, and will see
energy created at a rate \(3H\) per existing mass plus enthalpy. The present
density of the universe is roughly of the order
\(10^{-6}\,\mathrm{nucleons\,cm^{-3}}\), so with
\(H^{-1}=10^{10}\,\mathrm{yr}\) this would require an average creation of
order \(10^{-16}\,\mathrm{nucleons\,cm^{-3}\,yr^{-1}}\). The steady state
model is silent as to whether this new matter is created as hydrogen, or
protons plus electrons, or neutrons, and it does not tell us whether this
new matter appears near the old matter or in the depths of intergalactic
space. However, violent events do seem to be occurring in the nuclei of many
galaxies, so galactic nuclei seem like natural candidates for the location
of continuous creation.

The steady state cosmology makes very definite predictions as to the
correlation of luminosity distance with red shift. According to
Eq.~\eqref{eq:14.3.1}, if light leaves a comoving source at time \(t_1\) and
arrives at the origin at time \(t_0\), the source must be at a coordinate
\(r_1\) given for \(k=0\) by
\begin{equation}
  r_1
  =
  r(t_1)
  \equiv
  \int_{t_1}^{t_0}\frac{\dd t}{a(t)}
  =
  H^{-1}a^{-1}(t_0)
  \left\{
    \exp[H(t_0-t_1)]-1
  \right\}.
  \tag{14.8.9}\label{eq:14.8.9}
\end{equation}
Also, Eq.~\eqref{eq:14.3.6} gives the red shift of such a source as
\begin{equation}
  z=\exp[H(t_0-t_1)]-1,
  \tag{14.8.10}\label{eq:14.8.10}
\end{equation}
so the luminosity distance of the source is given by
Eq.~\eqref{eq:14.4.14} as
\begin{equation}
  d_L(z)=H^{-1}z(1+z).
  \tag{14.8.11}\label{eq:14.8.11}
\end{equation}
As a check, note that Eqs.~\eqref{eq:14.8.9}--\eqref{eq:14.8.11} agree
with Eqs.~\eqref{eq:14.6.6}, \eqref{eq:14.6.4}, and
\eqref{eq:14.6.8} for the appropriate deceleration parameter \(q=-1\).
This value does not seem to agree with the \(q_0\) determined from the
observed \(d_L\)-versus-\(z\) relation, as discussed in
Section~\ref{sec:14.6}.

The ``angular diameter'' distance \(d_A\) is given in this model by
Eqs.~\eqref{eq:14.4.22} and \eqref{eq:14.8.11} as
\begin{equation}
  d_A(z)=\frac{H^{-1}z}{1+z}.
  \tag{14.8.12}\label{eq:14.8.12}
\end{equation}
Note that \(d_A(z)\) approaches the finite constant \(H^{-1}\) as
\(z\to\infty\). Hence objects with large red shifts look very faint, but
their angular diameters do not shrink below a minimum value. If \(H^{-1}\)
is \(3\times10^9\,\mathrm{pc}\), then a galaxy of diameter
\(10^4\,\mathrm{pc}\) will never appear smaller than about \(0.6''\).

If we count the number of sources with red shift less than \(z\), we must
look back to a time \(t_z\) given by Eq.~\eqref{eq:14.7.8} as
\begin{equation}
  t_z=t_0-H^{-1}\ln(1+z).
  \tag{14.8.13}\label{eq:14.8.13}
\end{equation}
To count the number of sources with apparent luminosity greater than \(l\),
we must look back to a time given by Eq.~\eqref{eq:14.7.9}, which together
with Eq.~\eqref{eq:14.8.9} now reads
\[
  \exp[H(t_0-t_l)]
  \left\{
    \exp[H(t_0-t_l)]-1
  \right\}
  =
  \left(
    \frac{LH^2}{4\pi l}
  \right)^{1/2}.
\]
The solution is
\begin{equation}
  t_l(L)
  =
  t_0
  -
  H^{-1}
  \ln
  \left[
    \frac12
    +
    \left\{
      \frac14
      +
      \left(
        \frac{LH^2}{4\pi l}
      \right)^{1/2}
    \right\}^{1/2}
  \right].
  \tag{14.8.14}\label{eq:14.8.14}
\end{equation}
The \(t_1\)-integral in Eq.~\eqref{eq:14.7.7} can be done explicitly, and
we find the number of sources with red shift less than \(z\) and apparent
luminosity greater than \(l\) as
\begin{equation}
  N(<z,>l)
  =
  \int_0^\infty
  n(L)
  \min\!\left\{
    V(t_z),V(t_l(L))
  \right\}
  \dd L,
  \tag{14.8.15}\label{eq:14.8.15}
\end{equation}
where \(V\) is the volume
\begin{align}
  V(t)
  &\equiv
  \int_t^{t_0}
  4\pi r^2(t_1)a^2(t_1)\,\dd t_1
  \notag\\
  &=
  4\pi H^{-3}
  \left\{
    H(t_0-t)
    -\frac32
    +2\exp[-H(t_0-t)]
    -\frac12\exp[-2H(t_0-t)]
  \right\},
  \tag{14.8.16}\label{eq:14.8.16}
\end{align}
and \(n(L)\dd L\) is the time-independent proper number density of sources
with absolute luminosity between \(L\) and \(L+\dd L\).

As a special case of Eq.~\eqref{eq:14.8.15}, we note that the number of
sources with red shift less than \(z\) is
\begin{equation}
  N(<z)
  =
  4\pi H^{-3}n
  \left\{
    \ln(1+z)
    -
    \frac{z(1+3z/2)}{(1+z)^2}
  \right\},
  \tag{14.8.17}\label{eq:14.8.17}
\end{equation}
where \(n\) is the total number density of sources,
\[
  n\equiv\int_0^\infty n(L)\,\dd L.
\]
This result is independent of any assumptions concerning the luminosity
distribution of sources, and of course no evolution of the source density
or luminosity distribution is possible in a steady state universe. However,
with the limited statistics now available, it appears that
Eq.~\eqref{eq:14.8.17} does not agree with the observed red-shift
distribution of the quasi-stellar radio
sources.\hyperref[ch14-ref-67]{\textsuperscript{67}} In particular, the
observed red-shift distribution of the quasi-stellar sources shows a
pronounced peak\hyperref[ch14-ref-57]{\textsuperscript{57}} near \(z=1.95\),
which is absent in Eq.~\eqref{eq:14.8.17}. It should be noted, though, that
the observed value of \(N(<z)\) should in general be smaller than the
theoretical prediction \eqref{eq:14.8.17}, because some sources are not
counted if their optical or radio strength is too small.

As another special case of Eq.~\eqref{eq:14.8.15}, we note that the number
of sources with apparent luminosity greater than \(l\) is conveniently
written by defining
\[
  X(L,l)
  \equiv
  \frac12
  +
  \left[
    \frac14
    +
    \left(
      \frac{LH^2}{4\pi l}
    \right)^{1/2}
  \right]^{1/2}.
\]
Then
\begin{equation}
  N(>l)
  =
  4\pi H^{-3}
  \int_0^\infty
  n(L)
  \left\{
    \ln X(L,l)
    -\frac32
    +2X^{-1}(L,l)
    -\frac12X^{-2}(L,l)
  \right\}
  \dd L.
  \tag{14.8.18}\label{eq:14.8.18}
\end{equation}
In contrast with \(N(<z)\), the result here depends on the details of the
distribution function \(n(L)\).

A quantity of greater observational interest is the number of radio sources
with strength at frequency \(\nu\) greater than \(S\). If all sources have
the same ``straight'' spectrum \(P\propto\nu^{-\alpha}\), then
Eq.~\eqref{eq:14.7.15} gives the number of such sources as
\begin{equation}
  N(>S;\nu)
  =
  \int_0^\infty
  n(P,\nu)V(t_{S\alpha}(P))\,\dd P,
  \tag{14.8.19}\label{eq:14.8.19}
\end{equation}
where \(n(P,\nu)\dd P\) is the number of sources with intrinsic power at
frequency \(\nu\) between \(P\) and \(P+\dd P\);
\(V(t)\) is given by Eq.~\eqref{eq:14.8.16}; and \(t_{S\alpha}(P)\) is
determined by Eq.~\eqref{eq:14.7.16}, which now reads
\begin{align}
  &\exp\!\left[
    \frac12(3+\alpha)H(t_0-t_{S\alpha})
  \right]
  -
  \exp\!\left[
    \frac12(1+\alpha)H(t_0-t_{S\alpha})
  \right]
  \notag\\[-2pt]
  &\hspace{7.2cm}
  =
  \left(\frac{PH^2}{S}\right)^{1/2}.
  \tag{14.8.20}\label{eq:14.8.20}
\end{align}
This equation cannot be solved analytically for the observed spectral index
\(\alpha\simeq0.7\). However, it follows from
Eqs.~\eqref{eq:14.8.20}, \eqref{eq:14.8.19}, and \eqref{eq:14.8.16} that
\(N(>S,\nu)\) decreases more slowly than \(S^{-3/2}\) for all source
strengths, in contradiction with observed number counts, which decrease
more rapidly\hyperref[ch14-ref-63]{\textsuperscript{63}} than \(S^{-3/2}\)
for \(S\) greater than about
\(4\times10^{-26}\,\mathrm{W\,m^{-2}\,Hz^{-1}}\), and only then begin to
decrease more slowly than \(S^{-3/2}\). In the last section we saw that
these observations are also inconsistent with the results of non-steady-
state cosmologies, but in that case the discrepancy could be removed by
assuming an evolution of source densities, while in the steady state
cosmology no evolution of the source density is allowed.

As a check, we note that the quantities \(\beta_0(L)\) and
\(\beta_0(P,\nu)\) defined by Eqs.~\eqref{eq:14.7.22} and
\eqref{eq:14.7.23} must vanish in the steady state model:
\[
  \beta_0(L)=\beta_0(P,\nu)=0.
\]
Hence Eqs.~\eqref{eq:14.7.27}, \eqref{eq:14.7.28}, and
\eqref{eq:14.7.29}, with \(q_0=-1\) and straight spectra, now give the
number counts for ``nearby'' sources as
\begin{equation}
  N(<z)
  =
  \frac{4\pi}{3}H^{-3}z^3n
  \left\{
    1-\frac94z+\cdots
  \right\},
  \tag{14.8.21}\label{eq:14.8.21}
\end{equation}
\begin{equation}
  N(>l)
  =
  \frac{4\pi}{3}(4\pi l)^{-3/2}
  \int_0^\infty
  n(L)
  \left\{
    1
    -\frac{21}{4}
    \left(\frac{LH^2}{4\pi l}\right)^{1/2}
    +\cdots
  \right\}
  L^{3/2}\,\dd L,
  \tag{14.8.22}\label{eq:14.8.22}
\end{equation}
\begin{align}
  N(>S,\nu)
  =
  \frac{4\pi}{3}S^{-3/2}
  \int_0^\infty
  n(P,\nu)
  \left\{
    1
    -\frac34(2\alpha+5)
    \left(\frac{PH^2}{S}\right)^{1/2}
    +\cdots
  \right\}
  P^{3/2}\,\dd P.
  \tag{14.8.23}\label{eq:14.8.23}
\end{align}
These agree with the power-series expansions of the general formulas
\eqref{eq:14.8.17}, \eqref{eq:14.8.18}, and \eqref{eq:14.8.19}.

The steady state model does not appear to agree with the observed
\(d_L\)-versus-\(z\) relation or with the source counts for \(N(<z)\) and
\(N(>S,\nu)\). In a sense, this disagreement is a credit to the model; alone
among all cosmologies, the steady state model makes such definite predictions
that it can be disproved even with the limited observational evidence at our
disposal. The steady state model is so attractive that many of its adherents
still retain hope that the evidence against it will eventually disappear as
observations improve. However, if the cosmic microwave radiation discussed
in the next chapter is really black-body radiation, it will be difficult to
doubt that the universe has evolved from a hotter, denser early stage.
